% TOP_PROBLEM: {"id":30007517,"problem_number":"LOCAL-30007517","title":"Recession hysteresis","category_id":21,"category":{"id":21,"name":"economics","display_name":"Economics","description":"Open economic theory statements, published research agendas, and proposed scoped specifications; question provenance and historical priority are qualified per record.","slug":"economics","order_index":21,"created_at":"2026-10-08T00:00:00Z"},"status":"research_question","external_url":null,"legacy_ids":[],"published":false,"statement":"Identify persistent productive-capacity losses from a temporary demand contraction and the extent to which a defined stabilization intervention reverses them.","economics":{"original_id":6,"field":"Business cycles and macro dynamics","kind":"identification","formalization_basis":"adapted_research_question","source_status":"explicit_open","priority_status":"earliest_found","priority_note":"This is the earliest explicit relevant statement verified in this audit, not a claim of original priority; older literature and earlier versions may contain antecedents.","earliest_verified_reference":{"authors":["Janet L. Yellen"],"title":"Macroeconomic Research After the Crisis","year":2016,"url":"https://www.federalreserve.gov/newsevents/speech/yellen20161014a.htm","doi":null,"locator":"Section “The Influence of Demand on Aggregate Supply,” paragraphs on reversing hysteresis and research needs","evidence_summary":"The speech explicitly asks whether a high-pressure economy can reverse adverse supply effects of weak demand and calls for research on the size, benefits, and costs of this mechanism."},"current_reference":{"authors":["Janet L. Yellen"],"title":"Macroeconomic Research After the Crisis","year":2016,"url":"https://www.federalreserve.gov/newsevents/speech/yellen20161014a.htm","doi":null,"locator":"Section “The Influence of Demand on Aggregate Supply,” paragraphs on reversing hysteresis and research needs","evidence_summary":"The speech explicitly asks whether a high-pressure economy can reverse adverse supply effects of weak demand and calls for research on the size, benefits, and costs of this mechanism."},"formalization_note":"This translates the speech’s explicit reversal question into potential-outcome estimands; permanent limits require additional assumptions."},"statusQualification":"Published question or research agenda verified at the cited locator. The mathematical specification is adapted for this catalogue and includes additional assumptions; source publication does not establish first-ever priority or certify this exact model remains unsolved. This translates the speech’s explicit reversal question into potential-outcome estimands; permanent limits require additional assumptions.","statusReviewedAt":"2026-10-08"}
% FIRST_POSTED: 2026-10-08
% ENTRY_KIND: statement_only
% !TeX program = lualatex
\documentclass[11pt]{article}
\usepackage{fontspec}
\usepackage[a4paper,margin=25mm]{geometry}
\usepackage{amsmath,amssymb,mathtools}
\usepackage{xurl}
\usepackage[hidelinks]{hyperref}
\newcommand{\sref}[2]{\hyperlink{source:#1:#2}{[S#2]}}
\setlength{\emergencystretch}{3em}
\begin{document}
% problemId:30007517
\section*{Recession hysteresis}\label{30007517}
\subsection{Definitions and mathematical statement}
Consider regions or countries \(i\) observed annually, with potential-output index \(K_{it}\) defined before estimation from productive capital, labor capacity, and measured technology. Let \(d\in\{0,1\}\) denote exposure to an externally identified temporary demand contraction, and let \(a\in\{0,1\}\) denote a subsequently implemented stabilization intervention. The intervention must be defined by its instruments, duration, and financing, rather than by realized output growth. Write \(K_{i,t+h}(d,a)\) for the potential outcome under each joint regime.
For horizons after the demand disturbance has ended, define
\[
\tau_h=E[\log K_{i,t+h}(1,0)-\log K_{i,t+h}(0,0)],\quad
\rho_h=E[\log K_{i,t+h}(1,1)-\log K_{i,t+h}(1,0)].
\]
Identify whether the contraction causes a persistent capacity loss and how much the intervention reverses it. Report finite-horizon estimates and the assumptions needed to extrapolate \(\lim_{h\to\infty}\tau_h\) or \(\rho_h\); finite persistence alone cannot identify permanence. Separate demand exposure from correlated supply destruction, policy selection, migration, and measurement revisions. Quantify mechanisms through investment, labor-market attachment, or innovation using independently justified mediation assumptions. Compare the discounted capacity gain with intervention costs under stated welfare weights. The target is scoped causal hysteresis and reversibility, not a claim that every recession permanently lowers output.

\par\textbf{Formulation and scope.} This translates the speech’s explicit reversal question into potential-outcome estimands; permanent limits require additional assumptions.

\par\textbf{Open-question provenance.} An explicit question or research agenda was verified at the source locator. This does not by itself certify that every later version remains unresolved \sref{30007517}{1}.

\par\textbf{Historical priority.} This is the earliest explicit relevant statement verified in this audit, not a claim of original priority; older literature and earlier versions may contain antecedents.

\subsection{Short English statement}
Identify persistent productive-capacity losses from a temporary demand contraction and the extent to which a defined stabilization intervention reverses them.

\subsection{Sources}
\begin{itemize}\raggedright
\item[\textnormal{[S1]}]\hypertarget{source:30007517:1}{} Janet L. Yellen. 2016. Macroeconomic Research After the Crisis. Section “The Influence of Demand on Aggregate Supply,” paragraphs on reversing hysteresis and research needs.
\url{https://www.federalreserve.gov/newsevents/speech/yellen20161014a.htm}.
{\small\textit{Earliest verified question/agenda reference; historical priority qualified below. The speech explicitly asks whether a high-pressure economy can reverse adverse supply effects of weak demand and calls for research on the size, benefits, and costs of this mechanism.}}
\end{itemize}
\end{document}
